\section{Gravity}

Gravity emerges from the geometry. It is curvature-sourced and bulk-mediated, and it comes out with the right
law. There are two distinct gravitational objects here, the Newtonian potential on the flat cusp and the
holographic correlator in the curved bulk, and they should not be conflated. Beyond them the substrate reaches
the full apparatus of general relativity at the effective level, with the one exception noted at the end.

\subsection{The flat three-dimensional Newtonian $1/r$}

In the three-dimensional cusp the static potential of a point source is the lattice Green's function of the
Laplacian. Measured by the difference method, comparing $G(r)$ against $G(r_0)$ with the regular cosine
integrand, which cancels the $k = 0$ zero-mode contamination that wrecked naive finite-patch sums, the result is
a clean $1/r$, of the form $G(r) = a/r + b$ with $a \approx 0.083 \approx 1/(4\pi)$. This is the Newtonian
potential, and it is $1/r$ because the cusp is three-dimensional. The force is dimension-resolved.

\begin{result}[The Newtonian $1/r$ potential]
The static lattice Green's function in the three-dimensional cusp, measured by the difference method
\cite{pollard2026vibetest}, is $G(r) = a/r + b$ with $a \approx 0.083 \approx 1/(4\pi)$. The Newtonian potential
goes as $1/r$, fixed by the cusp being three-dimensional.
\end{result}

\subsection{The bulk-mediated holographic correlator}

A word on the term. \emph{Holography} is the idea that everything happening in a chunk of gravitating space can
be written on its lower-dimensional boundary, the way a flat hologram stores a three-dimensional scene
\cite{maldacena1998}. A correlator measures how strongly two points feel each other.

Between two boundary points the bulk-mediated correlator decays as a clean $1/r^2$. It was computed two ways. The
first is the exact Bethe-lattice cavity recursion, which gives $\mu = 1/b$ universally and exponent $\alpha = 2$,
validated against a directly-solved finite tree. The second is a tree-tensor contraction in the manner of the
multiscale entanglement renormalization ansatz \cite{ryu2006}, which gives roughly $1/r^{1.93}$. This is the
holographic two-point function, a different object from the flat-layer Newtonian $1/r$.

\begin{result}[The holographic $1/r^2$ correlator]
The bulk-mediated boundary-to-boundary correlator decays as $1/r^2$. The exact Bethe-lattice cavity recursion
gives exponent $\alpha = 2$ with $\mu = 1/b$, validated against a finite tree, and a tree-tensor contraction
gives $1/r^{1.93}$ \cite{pollard2026vibetest}. This is the holographic two-point function, distinct from the
Newtonian $1/r$.
\end{result}

\subsection{Full general relativity, not only Newton}

The substrate reaches full general relativity, to the level any emergent-gravity program reaches, not only the
Newtonian limit. Four pieces make the case.

The first is the Einstein field equations along Jacobson's entropic route \cite{jacobson1995}. Given the
substrate's measured area law $S = A/4G$ and the Unruh temperature $T = \kappa/2\pi$, the Clausius relation
$\delta Q = T\, dS$ on every local Rindler horizon is equivalent to the full nonlinear
$G_{\mu\nu} + \Lambda g_{\mu\nu} = 8\pi G\, T_{\mu\nu}$. The area-law coefficient fixes the Einstein coefficient
to exactly $8\pi$, with the verified value $2\pi/\eta = 25.133 = 8\pi$, the same $G$ as Newton's law, and the
weak field gives Poisson's equation with the correct $4\pi$ \cite{padmanabhan2010, verlinde2011}. So the area
law, a measured substrate fact, is the Einstein equation read as an equation of state.

\begin{result}[The Einstein equation as an equation of state]
From the measured area law $S = A/4G$ and the Unruh temperature $T = \kappa/2\pi$, the Clausius relation on every
local Rindler horizon is equivalent to $G_{\mu\nu} + \Lambda g_{\mu\nu} = 8\pi G\, T_{\mu\nu}$. The Einstein
coefficient is fixed to $8\pi$, verified as $2\pi/\eta = 25.133 = 8\pi$ \cite{pollard2026vibetest}, with the same
$G$ as the Newtonian law and the correct $4\pi$ in the weak-field Poisson equation.
\end{result}

The second is light bending. Because the metric curves time and space together, a photon deflects by
$4GM/b$, twice the naive Newtonian amount, with the verified ratio $2.0000$. This is the cleanest observable
separating full general relativity from Newton, and the substrate gets it right.

\begin{result}[Light bending equals $4GM/b$]
A photon passing a mass $M$ at impact parameter $b$ deflects by $4GM/b$, twice the Newtonian value, with measured
ratio $2.0000$ \cite{pollard2026vibetest}. The factor of two is the signature of full general relativity, and the
substrate reproduces it.
\end{result}

The third is black-hole thermodynamics, complete and self-consistent. From $S = A/4$ and $T = 1/(8\pi M)$ follow
the first law $dM = T\, dS$ with error around $10^{-9}$, the Smarr formula $M = 2TS$ exactly, the Bekenstein
bound saturated, a negative heat capacity $C = -8\pi M^2$ which is the evaporation instability, and a lifetime
scaling as $M^3$. The substrate's own de Sitter horizon carries the same thermodynamics, with numbers fixed by
the measured growth, giving $T = H/2\pi$ and $\Lambda = 3H^2$.

The fourth is gravitational waves. A massless spin-two graviton lives on the flat cusp, with dispersion
$\omega = |k|$, $z = 1$, and exactly two transverse-traceless polarizations. A binary radiates the quadrupole
waveform, with gravitational-wave frequency twice the orbital frequency and power
$P = (32/5)\,\mu^2 a^4 \omega^6$. The inspiral chirp follows $f \sim (t_c - t)^{-3/8}$, the LIGO exponent, with
the verified value $-0.375$.

\subsection{Why curvature is essential}

A flat lattice, the Euclidean $\{3,4,3,3\}$ with the same $D_4$ coin, has no curvature. It therefore has no
curvature-sourced gravity, no holographic boundary, and no expansion. This is why the substrate must be
hyperbolic. The flat twin is a clean control that isolates curvature as the source of gravity, set against the
hyperbolic case where the best of each dimension is compared.


Gravity is not added by hand. The flat three-dimensional cusp has a $1/r$ Newtonian potential, the curved bulk
mediates a clean $1/r^2$ holographic correlator, and the full general-relativistic relations follow at the
effective level, all falling out of the geometry.

\subsection{Caveat, the open frontier}

The naive finite-patch Green's function is unreliable, spoiled by the $k = 0$ mode and boundary contamination,
and the difference method is the correct tool that supersedes the earlier confounded measurements. The one thing
not done, and not done by anyone working from a micro-rule, is full nonlinear numerical relativity simulated cell
by cell, a binary merger or a forming horizon evolved from the ternary tone directly, which not even lattice QCD
attempts. What the substrate has is the effective theory together with the quantitative relations, the Einstein
equation of state, the limits, the thermodynamics, and the waveforms, each tied to substrate-measured inputs, the
area law, the light cone, the flat cusp, and the growth rate $H$. That is what it precisely means to say the
substrate supports general relativity.
